% !TEX root = ../main.tex
\section{Related Work}
The decomposition of nonstationary signals into time-varying frequency components is classically addressed by the Hilbert--Huang Transform (HHT), which extracts intrinsic mode functions (IMFs) and instantaneous frequencies \cite{huang1998empirical}. Unlike Fourier methods, HHT allows adaptive, nonstationary decomposition.
Recent work extends spectral modeling into neural frameworks. Neural spectral density estimation demonstrates that neural networks can approximate full spectral operators under mild conditions \cite{mohammadi2026spectral}. This provides theoretical grounding for learning frequency structure directly from data.
However, spectral entanglement remains a challenge: trends, cycles, and noise overlap in frequency space, particularly in nonstationary settings \cite{an2026fredn}. This motivates architectures capable of adaptive decomposition.
NS-SDN can be viewed as a neural generalization of IMF decomposition, where spectral components are learned endogenously rather than extracted via signal processing.